\chapter{光学仪器的基本原理}

\problem{4.1}{眼睛的构造可简化为一折射球面模型，其曲率半径为 $5.55\,\mathrm{mm}$，内部为折射率等于 $4/3$ 的液体，外部是空气，其折射率近似地等于 $1$。试计算眼球的两个焦距。用肉眼观察月球时，月球对眼的张角为 $1^\circ$，问视网膜上月球的像有多大？}
\answerarea

\problem{4.2}{把人眼的晶状体简化成距视网膜 $2\,\mathrm{cm}$ 的一个凸透镜，有人能看清距离在 $100\,\mathrm{cm}$ 到 $300\,\mathrm{cm}$ 间的物体。试问：（1）此人看清远点和近点时，眼睛透镜的焦距是多少？（2）为看清 $25\,\mathrm{cm}$ 远的物体，需配戴怎样的眼镜？}
\answerarea
\clearpage

\problem{4.3}{一照相机对准远物时，底片距物镜 $18\,\mathrm{cm}$，当镜头拉至最大长度时，底片与物镜相距 $20\,\mathrm{cm}$，求目的物在镜头前的最近距离？}
\answerarea

\problem{4.4}{两星所成的视角为 $4'$，用望远镜物镜照相，所得两像点相距 $1\,\mathrm{mm}$，问望远镜物镜的焦距是多少？}
\answerarea
\clearpage

\problem{4.5}{一显微镜具有三个物镜和两个目镜，三个物镜的焦距分别为 $16\,\mathrm{mm}$、$4\,\mathrm{mm}$ 和 $1.9\,\mathrm{mm}$，两个目镜的标示分别为 $5\times$ 和 $10\times$。设三个物镜造成的像都能落在像距为 $160\,\mathrm{mm}$ 处，问这显微镜的最大和最小的放大本领各为多少？}
\answerarea

\problem{4.6}{一显微镜物镜焦距为 $0.5\,\mathrm{cm}$，目镜焦距为 $2\,\mathrm{cm}$，两镜间距为 $22\,\mathrm{cm}$。观察者看到的像在无穷远处，试求物体到物镜的距离和显微镜的放大本领。}
\answerarea
\clearpage

\problem{4.7}{试证明望远镜的物镜为有效光阑（假定物镜和目镜的孔径相差不太悬殊）。}
\answerarea

\problem{4.8}{已知望远镜物镜的边缘即为有效光阑，试计算并作图求入射光瞳和出射光瞳的位置。}
\answerarea
\clearpage

\problem{4.9}{组成题 4.9 图的简单望远镜中各薄透镜的参数为 $\mathrm{L_1}$：$f_1=10\,\mathrm{cm}$，$D_1=4\,\mathrm{cm}$；$\mathrm{L_2}$：$f_2=2\,\mathrm{cm}$，$D_2=1.2\,\mathrm{cm}$；$\mathrm{L_3}$：$f_3=2\,\mathrm{cm}$，$D_3=1.2\,\mathrm{cm}$。计算该系统出射光瞳的位置和大小。\fig[0.54]{4_9}}
\answerarea

\problem{4.10}{有一光阑直径为 $5\,\mathrm{cm}$，放置在薄透镜后 $3\,\mathrm{cm}$ 处。透镜的焦距为 $5\,\mathrm{cm}$，孔径为 $6\,\mathrm{cm}$。现有一高为 $3\,\mathrm{cm}$ 的物 $PQ$ 置于透镜前 $12\,\mathrm{cm}$ 处。试求：（1）计算对主轴上 $P$ 点的入射光瞳和出射光瞳的大小和位置；（2）找到像的位置；（3）作光路图。}
\answerarea
\clearpage

\problem{4.11}{题 4.11 图中的 $H$、$H'$ 为光具组的主点，$F$、$F'$ 为焦点，$E$ 为对于物点 $P$ 的入射光瞳，$EO$ 为其半径。已知 $EO=2\,\mathrm{cm}$，$HP=20\,\mathrm{cm}$，$HF=15\,\mathrm{cm}$，$HO=5\,\mathrm{cm}$，$H'F'=15\,\mathrm{cm}$，$HH'=5\,\mathrm{cm}$，物长 $PQ=0.5\,\mathrm{cm}$。试作光路图并计算：（1）像的位置；（2）像长；（3）入射孔径角；（4）$P$ 点的出射光瞳半径和出射孔径角。\fig[0.52]{4_11}}
\answerarea

\problem{4.12}{一可认为是点光源的节能灯悬在圆桌中央的上空，桌的半径为 $R$，为了使桌的边缘能得到最大的照度，灯应悬在离桌面中心多高处？}
\answerarea
\clearpage

\problem{4.13}{一焦距为 $20\,\mathrm{cm}$ 的薄透镜，放在发光强度为 $15\,\mathrm{cd}$ 的点光源之前 $30\,\mathrm{cm}$ 处，在透镜后面 $80\,\mathrm{cm}$ 处放一屏，在屏上得到明亮的圆斑。求不计透镜中光的吸收时，圆斑的中心照度。}
\answerarea

\problem{4.14}{一长为 $5\,\mathrm{mm}$ 的线状物体放在一数码相机镜头前 $50\,\mathrm{cm}$ 处，在 CCD 上形成的像长为 $1\,\mathrm{mm}$。若 CCD 后移 $1\,\mathrm{cm}$，则像的弥散斑宽度为 $1\,\mathrm{mm}$。试求数码相机镜头的 $F$ 数。}
\answerarea
\clearpage

\problem{4.15}{某种玻璃在靠近钠光的黄色双谱线（其波长分别为 $589\,\mathrm{nm}$ 和 $589.6\,\mathrm{nm}$）附近的色散率 $\mathrm{d}n/\mathrm{d}\lambda$ 为 $-360\,\mathrm{cm^{-1}}$，求由此种玻璃制成的能分辨钠光双谱线的三棱镜，底边宽度应不小于多少？}
\answerarea

\problem{4.16}{设计一块平面透射光栅，要求满足：（1）使波长 $600\,\mathrm{nm}$ 的第二级谱线的衍射角小于 $30^\circ$，并能分辨 $0.02\,\mathrm{nm}$ 的波长差；（2）色散尽可能大；（3）第三级谱线缺级。试求：（1）光栅的缝宽、缝数、光栅常量和总宽度；（2）用这块光栅总共能看到 $600\,\mathrm{nm}$ 的几条谱线？}
\answerarea
\clearpage

\problem{4.17}{若要求显微镜能分辨相距 $0.000\,375\,\mathrm{mm}$ 的两点，用波长为 $550\,\mathrm{nm}$ 的可见光照明。试求：（1）此显微镜物镜的数值孔径 $R_{\mathrm{N.A.}}$；（2）若要求此两点放大后的视角为 $2'$，则显微镜的放大本领是多少？}
\answerarea

\problem{4.18}{夜间自远处驶来轿车的两前灯相距 $1.5\,\mathrm{m}$。如将眼睛的瞳孔看成产生衍射的圆孔，试估计视力正常的人在多远处才能分辨出光源是两个灯。设眼睛瞳孔的直径为 $3\,\mathrm{mm}$，设光源发出的光的波长 $\lambda$ 为 $550\,\mathrm{nm}$。}
\answerarea
\clearpage

\problem{4.19}{用孔径分别为 $20\,\mathrm{cm}$ 和 $160\,\mathrm{cm}$ 的两种望远镜能否分辨清月球上直径为 $500\,\mathrm{m}$ 的环形山？已知月球与地面的距离为地球半径的 $60$ 倍，而地球半径约为 $6\,370\,\mathrm{km}$。设光源发出的光的波长 $\lambda$ 为 $550\,\mathrm{nm}$。}
\answerarea

\problem{4.20}{电子显微镜的孔径角 $2u=8^\circ$，电子束的波长为 $0.1\,\mathrm{nm}$，试求：（1）它的最小分辨距离；（2）若人眼能分辨在明视距离处相距 $6.7\times 10^{-2}\,\mathrm{mm}$ 的两点，则该显微镜的放大倍数是多少？}
\answerarea
\clearpage

\sproblem{4.21}{平行光垂直投射于宽度为 $4\,\mathrm{cm}$ 的理想平面透射光栅上，已知在衍射角为 $60^\circ$ 的方向上的角色散率为 $0.5\times 10^{-2}\,\mathrm{rad/nm}$，试求光栅在该方向上最大处的分辨本领 $P$？}
\answerarea

\sproblem{4.22}{$\mathrm{He\text{-}Ne}$ 激光器发出波长为 $632.8\,\mathrm{nm}$、截面直径为 $2\,\mathrm{mm}$ 的激光束投向月球。已知月球和地面的距离为 $3.76\times 10^5\,\mathrm{km}$。试问：（1）在月球上得到的光斑的直径有多大？（2）如果这个激光经扩束器扩束后截面直径分别为 $2\,\mathrm{m}$ 和 $5\,\mathrm{m}$，再发向月球，试问在月球表面上的光斑直径各为多大？}
\answerarea
\clearpage

\sproblem{4.23}{宇宙飞船中的一架照相机在离地面 $200\,\mathrm{km}$ 的太空拍摄地面上的物体，如果要求它能分辨地面上相距 $1\,\mathrm{m}$ 的两点，照相机镜头的口径至少要多大？设光的波长为 $550\,\mathrm{nm}$。}
\answerarea

\sproblem{4.24}{已知月地距离约为 $3.8\times 10^5\,\mathrm{km}$，用上海天文台的口径为 $1.56\,\mathrm{m}$ 天体测量望远镜能分辨月球表面上两点的最小距离为多少？设光的波长为 $555\,\mathrm{nm}$。}
\answerarea
\clearpage

\sproblem{4.25}{为了分辨第二级钠光谱的双线（$589\,\mathrm{nm}$、$589.6\,\mathrm{nm}$），长度为 $15\,\mathrm{cm}$ 的平面光栅的光栅常量应为多少？}
\answerarea

\sproblem{4.26}{研究性课题：以 CD 或 DVD 光盘作为反射光栅，设计一光谱仪。}
\answerarea
\clearpage

\sproblem{4.27}{试绘出房门中设置的“猫眼”的结构及其光路图。}
\answerarea
\clearpage
